\section{Architecture}

\subsection{High-Level Design}

The system consists of 3 main components arranged in a sequential pipeline:

\begin{enumerate}[noitemsep]
    \item \textbf{Data pipeline:} Fetches historical lap data via FastF1, applies compound and track status filters, engineers features, and normalizes per race.
    % \item \textbf{Sequence builder:} Constructs sliding-window input sequences (10 laps $\times$ 9 features) and delta lap time targets per driver stint.
    \item \textbf{GRU model (per compound):} Three independent GRU networks trained on sequences filtered by tire compound (SOFT, MEDIUM, HARD).
    \item \textbf{Strategy simulator:} At inference time, loads a 2024 race, generates worn and fresh tire predictions for each candidate pit lap, and identifies the first lap where pitting is net positive.
\end{enumerate}

\subsection{Data Pipeline}

\subsubsection{Data Source and Filtering}

Race data is fetched using FastF1 v3.8.2 for the 2022 (22 rounds), 2023 (22 rounds), and 2024 (24 rounds) seasons. The 2024 season is held out 
entirely for evaluation; models are trained exclusively on 2022--2023 data. The model is trained on a temporal split between seasons rather than a random split across all seasons in order to prevent data leakage and ensure evaluation integrity.
Races within the same season share car perofmrance and trajectories and ragulatory context, so random samplying would allow the model to learn season specific patterns that would inflate the accuracy results without reflecting generalization.
For example, if the model is trained on car data from some races in 2024, then it could use the learned pattern from 2024 to perfectly predict
 the pit stop timing for the rest of the 2024 races that reside in the evaulation set.

Three filters are applied to each race:
\begin{enumerate}[noitemsep]
    \item \texttt{IsAccurate == True}: removes laps with timing errors
    \item \texttt{TrackStatus == '1'}: retains only green-flag laps (removes safety car, yellow flag, and red flag laps)
    \item \texttt{PitInTime.isna()}: removes the pit-entry lap where the driver lifts off
\end{enumerate}
Only dry-compound laps (SOFT, MEDIUM, HARD) are retained, and wet condition tires such as `intermediates' and `wet' are excluded.
 Weather data (air temperature, track temperature) is merged onto each lap via \texttt{pd.merge\_asof} on lap start time.

\subsubsection{Feature Engineering}

The input feature vector for each lap consists of 9 features:

\begin{table}[h]
\centering
\caption{Input Features}
\begin{tabular}{lll}
\toprule
Feature & Description & Normalization \\
\midrule
LapTimeSeconds & Lap time (s) & Per-race MinMax \\
StintLength & Laps on current tyre set & Per-race MinMax \\
FuelLoad & Estimated fuel remaining (kg) & Per-race MinMax \\
AirTemp & Ambient air temperature (\textdegree C) & Per-race MinMax \\
TrackTemp & Track surface temperature (\textdegree C) & Per-race MinMax \\
Compound\_SOFT & 1 if on SOFT tyres & Binary \\
Compound\_MEDIUM & 1 if on MEDIUM tyres & Binary \\
Compound\_HARD & 1 if on HARD tyres & Binary \\
TrackTemp\_global & TrackTemp scaled across all training races & Global MinMax \\
\bottomrule
\end{tabular}
\end{table}

Fuel load is estimated as $\max(0, 110 - \text{LapNumber} \times 1.6)$ kg, using a standard burn rate of 1.6 kg/lap from a 110 kg fuel load.

Continuous features are then normalized to $[0, 1]$ on a per-race basis using scikit-learn's \texttt{MinMaxScaler}:
\begin{equation}
\tilde{x}_i^{(r)} = \frac{x_i^{(r)} - x_{\min}^{(r)}}{x_{\max}^{(r)} - x_{\min}^{(r)}}
\end{equation}
where $r$ indexes the race and $x_{\min}^{(r)}, x_{\max}^{(r)}$ are computed over all laps in that race. Per-race scaling accounts for the wide variation in baseline lap times across circuits (e.g., Monaco 75 s vs.\ Bahrain 97 s); without it, the model would have to first learn circuit identity before learning degradation. The trade-off is that absolute temperature differences across circuits are erased. The \texttt{TrackTemp\_global} feature is therefore added separately, normalized once across all training races rather than per race, to preserve that signal.

\subsubsection{Sequence Building}

For each driver stint, a sliding window of width 10 generates training samples. A window starting at lap $i$ produces input $X_i \in \mathbb{R}^{10 \times 9}$ and target $y_i = t_{i+10} - t_{i+9}$ (the lap-to-lap delta of the next lap). Stints shorter than 10 laps produce no samples and are discarded. After splitting by compound, the training dataset contains:

\begin{table}[h]
\centering
\caption{Training Sequences by Compound (2022--2023)}
\begin{tabular}{lrr}
\toprule
Compound & Train Sequences & Test Sequences (2024) \\
\midrule
SOFT   & 4,087  & 200 \\
MEDIUM & 9,853  & 442 \\
HARD   & 15,041 & 2,584 \\
\midrule
Total  & 28,981 & 3,226 \\
\bottomrule
\end{tabular}
\end{table}

\subsection{Model Architecture}

The project evalutes two primary recurrent architectures: TireLSTM and TireGRU. Both take a sequence of shape $(B, 10, 9)$ as input and produce 
a scalar delta prediction.

\textbf{TireGRU} (selected final architecture). At each timestep $t$, the GRU computes a reset gate $\mathbf{r}_t$, an update gate $\mathbf{z}_t$, a candidate hidden state $\tilde{\mathbf{h}}_t$, and the new hidden state $\mathbf{h}_t \in \mathbb{R}^{64}$:
\begin{align}
\mathbf{r}_t &= \sigma(W_r \mathbf{x}_t + U_r \mathbf{h}_{t-1} + \mathbf{b}_r) \\
\mathbf{z}_t &= \sigma(W_z \mathbf{x}_t + U_z \mathbf{h}_{t-1} + \mathbf{b}_z) \\
\tilde{\mathbf{h}}_t &= \tanh\!\left(W_h \mathbf{x}_t + U_h(\mathbf{r}_t \odot \mathbf{h}_{t-1}) + \mathbf{b}_h\right) \\
\mathbf{h}_t &= (1 - \mathbf{z}_t) \odot \mathbf{h}_{t-1} + \mathbf{z}_t \odot \tilde{\mathbf{h}}_t
\end{align}
where $\sigma$ is the logistic sigmoid and $\odot$ denotes element-wise multiplication. After 10 timesteps, the final hidden state is projected to a scalar prediction:
\begin{equation}
\hat{y} = \text{Linear}(\text{Dropout}(\mathbf{h}_{10})) \in \mathbb{R}
\end{equation}

\noindent Configuration: hidden size 64, 2 stacked layers, dropout 0.2, batch-first. Total parameters: 39{,}425.

The model is trained to minimize mean squared error on the lap-to-lap delta target:
\begin{equation}
\mathcal{L} = \frac{1}{N}\sum_{i=1}^{N}(\hat{y}_i - y_i)^2
\end{equation}
Optimization uses the Adam optimizer~\cite{kingma2015} with learning rate $10^{-3}$, ReduceLROnPlateau scheduler (patience=3, factor=0.5), gradient clipping at norm 1.0, and early stopping at patience=7 on validation loss.

An LSTM with a learned attention pooling layer was also evaluated. A linear projection scores each of the 10 LSTM hidden states, a softmax over the timestep axis produces attention weights, and the context vector is the weighted sum of the hidden states before the output projection:
\begin{align}
\alpha_t &= \mathrm{softmax}_t(\mathbf{w}^\top \mathbf{h}_t) \\
\mathbf{c} &= \sum_{t=1}^{10} \alpha_t\, \mathbf{h}_t, \quad \hat{y} = \text{Linear}(\text{Dropout}(\mathbf{c}))
\end{align}
This is a content-based attention pooling rather than the additive query--key formulation of~\cite{bahdanau2015}: there is no separate query vector and no $\tanh$ nonlinearity. This variant was the worst performer (MAE $-22.5\%$ vs.\ baseline), as discussed in Section 4.3.

\subsection{Strategy Simulator}

At inference time, the simulator loads a 2024 race for a specific driver, normalizes the data using the same per-race scaler, and evaluates each 
candidate pit lap as follows:

\textbf{Worn tire projection} (\texttt{predict\_worn}): Starting from the last 10 observed laps, the model auto-regressively predicts 20 future laps. Because the GRU outputs a delta $\Delta\hat{t}_i$ rather than an absolute lap time, predicted lap times are reconstructed by integration:
\begin{equation}
\hat{t}_{i+1} = \hat{t}_i + \Delta\hat{t}_i
\end{equation}
seeded from $\hat{t}_0$ equal to the last observed lap time. At each rollout step, StintLength increments by 1 and FuelLoad decrements by the per-lap burn rate; the compound one-hot remains constant.

\textbf{Fresh tire projection} (\texttt{predict\_fresh}): Same integration scheme, but seeded from the earliest low-tyre-life laps of the target compound (StintLength $< 3$ normalized laps) earlier in the race, with StintLength reset to 0 to simulate a fresh stint.

\textbf{Decision criterion}: For each candidate pit lap $\ell$, compute:
\begin{equation}
\Delta(\ell) = \underbrace{L_\text{pit} + \sum_{k=1}^{20} \hat{t}_k^\text{fresh}}_{\text{pit cost}} - \underbrace{\sum_{k=1}^{20} \hat{t}_k^\text{worn}}_{\text{stay cost}}
\end{equation}
where $L_\text{pit}$ is the normalized pit stop time loss (22 seconds / max lap time). Pit loss values are sourced per circuit rather than remaining constant. Las Vegas, Singapore, Monaco, and Miami are shown to be deviating notably from the 22 second default value.
Since $L_\text{pit}$ defines the threshold a stop must overcome to be net-positive, a uniform value would systematically over-recommend stops at slow pit-lane circuits and under-recommend them at fast ones, shifting the optimal pit window by several laps.
 The recommended pit lap is the first $\ell$ where $\Delta(\ell) < 0$. One-stop and two-stop strategies are evaluated by brute-force search over all valid $(p_1)$ and $(p_1, p_2)$ pit lap pairs; the total race time under each strategy determines which is preferred.
